%% =============================================================================
%% appendices/appendix-decisions.tex
%% Harmonia Architecture Decision Register
%% =============================================================================

\chapter{Architecture Decisions}
\label{app:decisions}

This appendix records architecture choices that constrain how Harmonia requirements are realised. Decisions are distinct from requirements: a requirement states what Harmonia must achieve; a decision records the architectural approach selected to achieve or protect that outcome.

\section{Current Decisions}

\subsection{ADR-001 --- Petasos Owns Messaging Runtime}
Petasos is the standalone Harmonia messaging capability and owns the ActiveMQ Artemis runtime. Ponos consumes Petasos messaging services and owns workflow/work execution; it does not own broker lifecycle.

\subsection{ADR-002 --- Pylai Owns External Interface and Protocol Behaviour}
Pylai owns external protocol/interface behaviour, including HL7/FHIR ingress and egress. Domain workflow logic remains outside the protocol boundary.

\subsection{ADR-003 --- Mnemosyne Owns Durable Application State}
Mnemosyne is the durable persistence capability. Mneme/Infinispan provides operational/cache state and must not redefine the existence of durable governed information.

\subsection{ADR-004 --- Calliope Owns Canonical Models and Semantic Governance}
Calliope owns canonical model definitions, profiles, terminology bindings, and semantic governance used by Harmonia processing.

\subsection{ADR-005 --- Themis Owns Security Policy Decisions}
Themis is the common Harmonia authorisation/policy capability. Authentication boundaries establish trusted identity; Themis evaluates permissions.

\subsection{ADR-006 --- FHIR R5 Is the Clinical Interoperability Boundary}
The durable clinical interoperability service is represented and exposed using HL7 FHIR R5. Proprietary source models are translated into the governed representation rather than propagated as the enterprise integration contract.

\subsection{ADR-007 --- Harmonia Is an Interoperability Authority, Not Universal Clinical Truth}
Source systems may remain authoritative for the activities they manage. Harmonia maintains the governed, vendor-neutral interoperability representation and preserves provenance and Information Authority rather than claiming universal clinical source-of-truth status.

\subsection{ADR-008 --- Iris-Clinical Begins as a Longitudinal Record Viewer}
The initial Iris-Clinical capability is read-oriented. The architecture permits selective clinical authoring later, but authoring is introduced only where terminology, validation, security, governance, and workflow prerequisites are satisfied.

\subsection{ADR-009 --- Information Authority Is Cross-Cutting}
AUTHORITATIVE, INFORMATIONAL, and ANECDOTAL classification is a core Harmonia concept rather than an Iris or Clinical-only feature.

\subsection{ADR-010 --- Clinical Search Must Be Authoritative-Backed}
Cache contents are not the definition of the clinical record. Clinical read/search must remain correct from durable state after cache loss or restart.

\subsection{ADR-011 --- Agora Uses Restricted Self-Hosted Matrix/Synapse}
The initial Agora collaboration environment is self-hosted, restricted, and non-federated. Patient collaboration rooms and membership are governed by Harmonia rather than unmanaged public federation.

\subsection{ADR-012 --- Paradeigma Is an Isolated Exemplar/Test Subsystem}
Paradeigma exists to exercise Harmonia interfaces and workflows with deterministic synthetic clinical scenarios. It is not a production clinical dependency and must remain isolated from production architecture.

\subsection{ADR-013 --- Kleio Owns Audit Evidence and Provenance}
Kleio is the Harmonia capability responsible for immutable audit evidence and governed provenance. Themis, Pylai, Ponos, Ergon, Iris, and other Harmonia subsystems may produce evidence, but do not independently own the enterprise audit model or durable audit record. Accepted audit evidence is append-only and shall not be updated, replaced, patched, deleted, or resurrected through ordinary application persistence mechanisms.

\subsection{ADR-014 --- Petasos Owns Durable Processing Transition Boundaries}
Petasos governs the durable transfer of work between independently recoverable Harmonia processing activities. Ponos owns execution and orchestration of those activities; Petasos establishes the durable boundary at which responsibility for work is transferred. ActiveMQ Artemis is an implementation mechanism used by Petasos and does not itself define Harmonia processing semantics.

\subsection{ADR-015 --- Replay Is Anchored to Explicit Durable Transitions}
Replay is supported only from Petasos-governed durable transition points explicitly designated as replayable. Harmonia does not provide arbitrary replay from internal implementation states. A deliberate replay creates a new processing attempt linked to the original transition through correlation, causation, and provenance; it does not modify, replace, or impersonate the original event or processing attempt. Replay, deterministic reprocessing, transport redelivery, and redelivery of an established result are distinct operations and shall not be treated as equivalent.

\subsection{ADR-016 --- Audit Is Anchored to Architecturally Significant Transitions}
Kleio AuditEvents are generated for architecturally significant actions and transitions, including acceptance or rejection of work, significant commitment or authority transitions, external dispatch, and deliberate replay where policy requires evidence. Internal implementation steps such as parsing, mapping helpers, individual validation rules, and ordinary method execution are not independently audited merely because they occur. Audit, provenance, and replay remain distinct concepts but may share a durable Petasos transition as their architectural anchor.

\subsection{ADR-017 --- Processing Pressure Is Expressed as Durable Backlog}
Where processing demand exceeds downstream capacity, Harmonia shall preferentially express pressure as bounded, observable, durable Petasos/Artemis backlog rather than unbounded growth in JVM heap, worker threads, database connections, or volatile work state. Queue depth, oldest-transition age, processing rate, retry/replay count, dead-letter count, and backlog drain rate form part of the operational model for transition processing.

\subsection{ADR-018 --- Mnemosyne Defines the Authoritative Durable State Boundary}
\label{adr:018-mnemosyne-durable-state-boundary}

\textbf{Status:} Accepted\\
\textbf{Date:} 24 September 2026

\subsubsection{Context}

Harmonia requires application state to remain available and recoverable across process, node and infrastructure failures while supporting high throughput, horizontal scaling and deployment across multiple failure domains.

The original Mneme architecture considered using Infinispan as more than a conventional cache. Under this model, application state would be synchronously replicated across multiple Infinispan instances, potentially distributed across independent sites or data centres. Each site could then independently persist replicated state through its local cache-store mechanism into a local Mnemosyne persistence service.

The intended benefits were:
\begin{itemize}
    \item high availability at the application and middleware tier rather than relying primarily upon database or storage-tier high availability;
    \item low write latency by allowing application processing to complete following distributed cache acceptance rather than waiting for physical database persistence;
    \item resilience to failure of an individual persistence store;
    \item independent persistence of replicated state at multiple sites;
    \item reduced dependence on expensive database replication and highly available storage infrastructure; and
    \item eventual reconciliation of persistence stores following temporary database, site or network failure.
\end{itemize}

Conceptually, the distributed Mneme fabric would therefore have formed the immediate acceptance boundary, with Mnemosyne persistence occurring asynchronously behind it.

This approach is technically feasible. However, for replicated cache state to constitute authoritative durable state, Harmonia would also need to guarantee recovery following cache, process, site and network failures. This would require Harmonia to own mechanisms for durable outstanding-write tracking, persistence retry, restart recovery, cross-site reconciliation, version consistency, conflict resolution, partition handling and eventual convergence of independent persistence stores.

The resulting capability would effectively introduce a distributed persistence and consistency layer above the existing durable persistence technologies.

Harmonia already provides distinct technologies appropriate to the principal resilience requirements:
\begin{itemize}
    \item \textbf{Mnemosyne/PostgreSQL} provides durable application-state persistence;
    \item \textbf{Petasos/Artemis} provides durable transfer of work between independently recoverable processing activities; and
    \item \textbf{Mneme/Infinispan} provides distributed caching and efficient access to reconstructable state.
\end{itemize}

The additional complexity required to make Mneme itself an authoritative persistence boundary is therefore not justified by the expected performance, availability or infrastructure-cost benefits.

\subsubsection{Decision}

\textbf{Mnemosyne defines the authoritative durable application-state boundary for Harmonia.}

Application state that requires durable persistence is not considered durably accepted solely because it has been written to, replicated by, or is currently available from Mneme/Infinispan.

A successful Mneme cache operation does not constitute durable application-state acceptance.

State requiring durable persistence is considered accepted when:
\begin{enumerate}
    \item it has been successfully committed through Mnemosyne to its authoritative durable persistence mechanism; or
    \item responsibility for further processing has been successfully transferred across an explicitly defined durable Petasos transition, where the state necessary to complete or reconstruct that processing is durably recoverable.
\end{enumerate}

Petasos/Artemis may therefore provide a durable \textbf{work acceptance} boundary, but does not replace Mnemosyne as the authoritative repository for committed application state.

Mneme/Infinispan will remain a distributed caching capability. State held by Mneme must be reconstructable from authoritative durable state or from an explicitly defined durable processing transition where reconstruction is part of that transition's contract.

Volatile or cached state must never be the sole recoverable representation of application state or work that Harmonia has acknowledged as durably accepted.

\subsubsection{Architectural Responsibilities}

\begin{center}
    \begin{tabular}{p{0.13\textwidth} p{0.29\textwidth} p{0.18\textwidth} p{0.29\textwidth}}
        \hline
        \textbf{Subsystem} & \textbf{Responsibility} & \textbf{Primary Technology} & \textbf{Durability Role} \\
        \hline
        Mneme & Distributed caching and efficient state access & Infinispan & Non-authoritative, reconstructable state \\
        Mnemosyne & Durable application-state persistence & PostgreSQL & Authoritative durable application state \\
        Petasos & Transfer of work between independently recoverable processing activities & ActiveMQ Artemis & Durable processing transitions \\
        Kleio & Immutable audit and provenance evidence & PostgreSQL & Durable immutable evidence \\
        \hline
    \end{tabular}
\end{center}

These responsibilities are deliberately distinct. Harmonia must not introduce implicit overlap between cache availability, processing durability and authoritative persistence.

\subsubsection{Consequences}

The decision simplifies Harmonia's failure and recovery semantics.

\paragraph{Positive Consequences}

\begin{itemize}
    \item Harmonia does not need to implement its own distributed database consistency and convergence mechanisms.
    \item Persistence failure can be represented explicitly rather than hidden behind successful cache operations.
    \item Database recovery and replication can use established PostgreSQL capabilities appropriate to each deployment topology.
    \item Petasos/Artemis can absorb processing pressure and temporary downstream outages as durable backlog rather than requiring volatile cache state to act as a persistence substitute.
    \item Mneme remains independently scalable and distributable without carrying an authoritative durability obligation.
    \item Cache contents may safely be evicted, restarted or reconstructed without causing loss of acknowledged authoritative state.
    \item Different deployment profiles may provide different levels of PostgreSQL, Artemis and Infinispan redundancy without changing Harmonia's logical durability semantics.
    \item Failure behaviour becomes easier to test and reason about.
\end{itemize}

\paragraph{Negative Consequences}

\begin{itemize}
    \item authoritative synchronous writes may incur greater latency than memory/cache acceptance;
    \item PostgreSQL remains part of the critical path for operations requiring immediate durable application-state commitment;
    \item deployments requiring high availability of authoritative state must provide an appropriate PostgreSQL HA/DR topology;
    \item Infinispan replication cannot by itself be used to mask loss of the authoritative persistence service; and
    \item some existing Harmonia persistence and cache interactions must be changed because they currently allow cache or JVM-local state to be treated as successfully persisted state.
\end{itemize}

These costs are accepted in preference to Harmonia assuming responsibility for distributed persistence, reconciliation and consistency.

\subsubsection{Failure Semantics}

The following invariant applies:

\begin{quote}
    \textbf{Accepted state must be durably recoverable.}
\end{quote}

This does \textbf{not} imply:

\begin{quote}
    Accepted state must always have already been committed to PostgreSQL.
\end{quote}

A durable Petasos transition may constitute acceptance of work where the information necessary to continue processing is durably recoverable from that transition.

However:

\begin{quote}
    \textbf{Cache presence, cache replication, JVM-local state or an asynchronous persistence attempt does not by itself constitute durable acceptance.}
\end{quote}

Consequently:
\begin{itemize}
    \item failure to persist required authoritative state must be propagated as failure;
    \item a local in-memory fallback must not convert unavailable durable persistence into apparent success;
    \item asynchronous processing may return acceptance only after the corresponding durable work boundary has been established;
    \item failure after a durable Petasos transition must result in recoverable backlog, redelivery or controlled failure rather than silent loss;
    \item cache loss must be recoverable from Mnemosyne or another explicitly defined durable source; and
    \item transient processing state such as a Pragma may remain volatile where it can be reconstructed without loss of acknowledged work.
\end{itemize}

\subsubsection{Rejected Alternative --- Mneme as a Distributed Persistence Fabric}

The rejected architecture would have treated replicated Infinispan state as Harmonia's primary immediate persistence/availability boundary.

Multiple independent Mneme instances or sites would replicate state between them. Each site would independently write replicated state through its cache-store implementation to a local Mnemosyne persistence service. Failure of an individual database or persistence site would not prevent application acceptance, with failed persistence operations subsequently retried and stores gradually reconciled.

This model was rejected \textbf{not because it is technically invalid}, but because completing it safely would require Harmonia to provide substantial additional distributed-systems capability, including:
\begin{itemize}
    \item durable tracking of outstanding persistence operations;
    \item recovery of pending writes following process or site restart;
    \item persistence retry and backoff;
    \item store convergence monitoring;
    \item version and ordering semantics;
    \item duplicate and idempotency handling;
    \item conflict detection and resolution;
    \item network-partition behaviour;
    \item reconciliation of independently persisted copies; and
    \item operational tooling for detecting and repairing divergence.
\end{itemize}

These responsibilities would duplicate capabilities already available through established database and durable-messaging technologies while materially increasing implementation, testing and operational complexity.

The likely reduction in persistence latency and infrastructure cost does not justify that complexity for Harmonia's expected workload and availability requirements.

\subsubsection{Relationship to Other Decisions}

This decision reinforces:
\begin{itemize}
    \item \textbf{ADR-003 --- Mnemosyne Owns Durable Application State}
    \item \textbf{ADR-010 --- Clinical Search Must Be Authoritative-Backed}
    \item \textbf{ADR-014 --- Petasos Owns Durable Processing Transition Boundaries}
    \item \textbf{ADR-015 --- Replay Is Anchored to Explicit Durable Transitions}
    \item \textbf{ADR-017 --- Processing Pressure Is Expressed as Durable Backlog}
\end{itemize}

It also establishes the architectural basis for the current implementation sequence:

\begin{verbatim}
Task 07   Remove volatile persistence fallback
             |
Task 08   Establish authoritative Clinical writes
             |
Task 09   Establish cache-aside point reads
             |
Task 10   Establish authoritative Clinical search
\end{verbatim}

Task 07 must remove cases where cache or JVM-local state is incorrectly treated as durable acceptance without prematurely redesigning the authoritative Clinical write path that is addressed by Task 08.

In alignment with ADR-020, Task 08 establishes authoritative CREATE and UPDATE write semantics, treating domain-appropriate lifecycle changes as conditional updates, without providing a physical authoritative DELETE path.

\subsubsection{Architectural Principle}

\begin{quote}
    \textbf{Petasos makes work recoverable.}\\
    \textbf{Mnemosyne makes application state durable.}\\
    \textbf{Mneme makes active state available, coordinated and fast.}\\
    \textbf{Kleio makes evidence permanent.}
\end{quote}

Harmonia will prefer these explicit responsibilities over introducing a distributed persistence protocol spanning them.

% =============================================================================
% ADR-019 --- Mneme Owns Distributed Resource Access and Coordination
% =============================================================================

\subsection{ADR-019 --- Mneme Owns Distributed Resource Access and Coordination}
\label{adr:019-mneme-distributed-resource-access-coordination}

\textbf{Status:} Accepted\\
\textbf{Date:} 25 September 2026

\subsubsection{Context}

Harmonia processes resources and data objects through validation, transformation, workflow, integration and persistence activities. Access to these objects is expected to be uneven, highly repetitive and strongly read-dominant.

At any point in time, only a comparatively small subset of the total persisted resource population is expected to be actively involved in processing. That active subset may be read repeatedly by different Harmonia components while it is validated, referenced, transformed, enriched, coordinated or progressed through workflows. A useful design expectation is approximately a 10:1 read-to-write ratio for this active working set; the ratio is indicative rather than a contractual limit.

Repeated authoritative retrieval from Mnemosyne for every such interaction would unnecessarily increase persistence traffic, database connection utilisation and processing latency.

Mneme was introduced to provide more than a convenient local cache. Using Infinispan and its distributed access mechanisms, Mneme is intended to provide a common low-latency resource-access and coordination capability across the Harmonia ecosystem.

The principal intended capabilities are:
\begin{enumerate}
    \item \textbf{Distributed resource availability} --- supported active resources and data objects can be accessed by Harmonia components regardless of which runtime node originally obtained, created or cached their working representation.
    \item \textbf{Distributed concurrency and version coordination} --- concurrent access to a given active resource or data object is coordinated through Mneme rather than independently by each client or JVM. Client code should not need to recreate distributed concurrency/version-management machinery.
    \item \textbf{Read-dominant performance} --- frequently read active state can be accessed at low latency without repeated database retrieval, while the comparatively smaller number of authoritative writes are committed through Mnemosyne.
    \item \textbf{Efficient internal distribution} --- native distributed-cache access and replication can be used within Harmonia without unnecessarily introducing HTTP network boundaries between internal components.
\end{enumerate}

The original Mneme design also contemplated \textbf{distributed durability}: replicated Infinispan state could act as an immediate durable acceptance fabric with persistence occurring behind it. ADR-018 deliberately rejects that responsibility because the required recovery, reconciliation, partition, convergence and durable outstanding-write semantics are not justified by the expected benefit.

Rejecting distributed durability does not reduce Mneme to an optional best-effort cache. Its distributed availability, coordination and performance semantics remain deliberate architectural responsibilities.

\subsubsection{Decision}

\begin{quote}
    \textbf{Mneme owns Harmonia's distributed resource access and coordination capability for active application state.}
\end{quote}

Mneme will use Infinispan to expose a distributed, low-latency and reconstructable representation of selected resources and data objects being actively used by Harmonia.

Mneme provides four defining characteristics:
\begin{itemize}
    \item \textbf{distributed availability} --- supported active state is accessible across Harmonia runtime instances through the distributed Mneme capability;
    \item \textbf{distributed coordination} --- Mneme provides the common coordination point for concurrent access and working versions of an individual active resource or data object;
    \item \textbf{read-dominant performance} --- Mneme absorbs high-volume, repetitive and potentially jittery reads against the active working set, reducing repeated access to authoritative persistence; and
    \item \textbf{reconstructability} --- Mneme state is non-authoritative and may be rebuilt from authoritative durable state or another explicitly defined durable source.
\end{itemize}

Mneme does \textbf{not} own distributed durability.

A successful Mneme operation, distributed replication, concurrency decision or working version does not by itself constitute authoritative durable acceptance. Authoritative application-state persistence and the final persisted version remain the responsibility of Mnemosyne.

\begin{quote}
    \textbf{Mneme coordinates active distributed state; Mnemosyne commits authoritative durable state.}
\end{quote}

The exact write-boundary protocol between Mneme coordination and Mnemosyne authoritative persistence is governed by the relevant write architecture and must preserve authoritative concurrency/version correctness at commit.

\subsubsection{Distributed Resource Availability}

Where a resource or data object is supported by Mneme, client code should be able to obtain its active representation through the Mneme capability without depending on the JVM or processing node that originally loaded or produced it.

Mneme therefore represents a system-wide resource-access capability, not a collection of unrelated per-process caches.

\begin{verbatim}
 Pylai      Ponos      Ergon      Praxis      other Harmonia components
   \          |          |          |                    /
    \         |          |          |                   /
     +--------+----------+----------+------------------+
                              |
                              v
                         +---------+
                         |  Mneme  |
                         |Infinispan|
                         +----+----+
                              |
                    authoritative state
                              |
                              v
                         +---------+
                         |Mnemosyne|
                         +---------+
\end{verbatim}

\subsubsection{Distributed Concurrency and Version Coordination}

Mneme is the distributed coordination point for concurrent access to an active resource or data object.

Where multiple Harmonia components or runtime nodes operate on the same active object, they must not independently establish incompatible local working versions merely because they execute in different JVMs.

Mneme should use appropriate Infinispan concurrency facilities so that the coordination mechanism is substantially transparent to ordinary client code. The implementation may use versioned or conditional operations, atomic replacement, locking or other appropriate Infinispan mechanisms, but those mechanisms are implementation choices beneath the Mneme contract.

This responsibility does not make Mneme the final authority for a durable write. Mnemosyne must still enforce the authoritative persistence transaction and persisted-version semantics.

\subsubsection{Read-Dominant Active Working Set}

Mneme is intentionally optimised around a comparatively small active subset of the total persisted resource population.

The expected access pattern is strongly read-dominant, with approximately ten reads for each write being a useful design assumption rather than a fixed service-level guarantee.

Typical processing may therefore involve repeated reads, validation, reference lookups, workflow evaluation and transformation against Mneme before a smaller number of create or update operations require authoritative persistence.

Mneme exists to prevent these repeated reads from degenerating into continuous database lookups while still allowing Mnemosyne to remain the authoritative durability boundary.

\subsubsection{Process-Local Fallback}

Where a component requires Mneme semantics:

\begin{quote}
    \textbf{Loss of Mneme must be visible as loss of distributed resource availability or coordination. It must not silently degrade into process-local application state.}
\end{quote}

Production services must therefore not use \texttt{ConcurrentHashMap}, local collections or equivalent process-memory structures as transparent substitutes for Mneme-backed application state.

Such a fallback would lose not only distributed caching, but also Mneme's resource-availability and concurrency/version-coordination semantics.

A component may respond to Mneme unavailability by failing the operation, or by using an explicitly designed degraded path where the architecture permits one. A degraded path must be deliberate and visible; it must not establish an alternative JVM-local Mneme.

This decision does \textbf{not} prohibit ordinary process-local memory. Local collections remain appropriate for temporary computation, immutable configuration, logger caches, local registrations, implementation metadata and other genuinely process-local concerns.

\begin{quote}
    \textbf{Process-local state is permitted. Process-local substitution for Mneme is not.}
\end{quote}

\subsubsection{Relationship Between Mneme and Mnemosyne}

For read-dominant access, the intended model is:

\begin{verbatim}
Request
   |
   v
 Mneme ---- HIT ----------------------> Return
   |
  MISS
   |
   v
Mnemosyne
   |
   +----> Populate/refresh Mneme
   |
   +----> Return
\end{verbatim}

For authoritative writes:

\begin{verbatim}
active distributed state / coordination
                 |
                 v
               Mneme
                 |
          authoritative write
                 |
                 v
             Mnemosyne
                 |
          durable commit/version
\end{verbatim}

The precise authoritative Clinical write path is addressed separately and must not be inferred solely from this diagram.

There is deliberately no transparent substitution of the form:

\begin{verbatim}
Mneme unavailable
       |
       v
ConcurrentHashMap
       |
       v
Pretend distributed semantics still exist
\end{verbatim}

\subsubsection{Consequences}

\paragraph{Positive Consequences}

\begin{itemize}
    \item supported active resources are available through a common distributed access capability across Harmonia runtime nodes;
    \item concurrency and working-version coordination are not independently reinvented by each client or process;
    \item high-volume and jittery reads can be absorbed without continuously querying Mnemosyne;
    \item native distributed-cache access can minimise unnecessary internal HTTP latency and protocol overhead;
    \item Mneme can be scaled and distributed independently of authoritative persistence;
    \item cache contents remain reconstructable and may be evicted or restarted without redefining durable truth; and
    \item the distinction between distributed coordination and durable authority is explicit.
\end{itemize}

\paragraph{Negative Consequences}

\begin{itemize}
    \item Mneme/Infinispan remains an important runtime dependency even though it is not a durability boundary;
    \item loss of Mneme may make distributed resource access or coordination unavailable until restored or an explicitly designed degraded path is used;
    \item concurrency/version behaviour must be designed and tested deliberately rather than assumed from cache replication alone;
    \item authoritative write paths must correctly bridge Mneme coordination and Mnemosyne persisted-version semantics; and
    \item appropriate monitoring of Mneme availability, topology and concurrency behaviour is operationally important.
\end{itemize}

These costs are accepted because distributed availability, coordination and read-dominant access are intentional Harmonia capabilities rather than incidental cache optimisations.

\subsubsection{Rejected Alternative --- Mneme as Distributed Durability}

ADR-018 records the rejected distributed-durability model. The idea remains technically attractive: replicated Infinispan state could form an immediate acceptance fabric and allow persistence stores to converge behind it.

It is rejected for Harmonia because doing it safely would require durable outstanding-write tracking, restart recovery, partition semantics, reconciliation, conflict resolution, convergence monitoring and associated operational tooling. Those costs are not justified by the expected workload and availability benefit.

This rejection applies specifically to \textbf{durability}. It does not reject Mneme's distributed availability, coordination, version-management or read-performance responsibilities.

\subsubsection{Rejected Alternative --- Process-Local Cache Fallback}

A JVM-local fallback may make an individual process appear available while silently destroying Mneme's distributed semantics.

\begin{verbatim}
Node A                 Node B                 Node C

Person/123 v7          Person/123 v6          Person/123 absent
Task/456 PROCESSING    Task/456 RECEIVED      Task/456 absent
\end{verbatim}

None of these copies need be authoritative for the situation to be incorrect. The nodes no longer share the resource availability or coordination model that Mneme exists to provide.

The apparent availability benefit therefore does not justify process-local substitution.

\subsubsection{Relationship to Other Decisions}

This decision complements:
\begin{itemize}
    \item \textbf{ADR-003 --- Mnemosyne Owns Durable Application State}
    \item \textbf{ADR-010 --- Clinical Search Must Be Authoritative-Backed}
    \item \textbf{ADR-014 --- Petasos Owns Durable Processing Transition Boundaries}
    \item \textbf{ADR-017 --- Processing Pressure Is Expressed as Durable Backlog}
    \item \textbf{ADR-018 --- Mnemosyne Defines the Authoritative Durable State Boundary}
\end{itemize}

Together:
\begin{itemize}
    \item \textbf{Mnemosyne} answers: Where does authoritative application state live and what version was durably committed?
    \item \textbf{Petasos} answers: Will accepted work survive until processing can continue?
    \item \textbf{Mneme} answers: How is active state made available, coordinated and read efficiently across Harmonia?
    \item \textbf{Kleio} answers: What durable evidence exists of significant actions and transitions?
\end{itemize}

\subsubsection{Architectural Principle}

\begin{quote}
    \textbf{Mneme makes active state available, coordinated and fast across Harmonia --- reconstructable, never authoritative, and never secretly local.}
\end{quote}

% =============================================================================
% ADR-020 --- Governed Information Uses Lifecycle State Rather Than Physical Deletion
% =============================================================================

\subsection{ADR-020 --- Governed Information Uses Lifecycle State Rather Than Physical Deletion}
\label{adr:020-governed-information-lifecycle-state}

\textbf{Status:} Accepted\\
\textbf{Date:} 25 September 2026

\subsubsection{Context}

In healthcare integration and interoperability environments, persisted resources (including patients, encounters, practitioners, clinical observations, documents, and processing tasks) represent clinical, legal, regulatory, and evidentiary facts.

Traditional persistence architectures often treat physical deletion (\texttt{DELETE} in SQL or HTTP \texttt{DELETE} in REST APIs) as a standard CRUD primitive that permanently removes records from durable storage. In a distributed health integration platform such as Harmonia, physical deletion introduces severe operational, clinical, and governance failures:

\begin{enumerate}
    \item \textbf{Loss of Clinical and Legal History} --- destroying persisted records erases the longitudinal record of care, compromises traceability, and violates healthcare data retention and evidentiary standards.
    \item \textbf{Referential Integrity Destruction} --- physical deletion of an entity breaks references in downstream systems, clinical documents, immutable audit records, and related resources across the enterprise.
    \item \textbf{Audit and Provenance Invalidation} --- accepted audit evidence governed by Kleio (ADR-013) must point to verifiable historical and provenance contexts; deleting the underlying entities destroys the audit graph.
    \item \textbf{Distributed State Ambiguity} --- in an ecosystem combining distributed caching (Mneme) and durable persistence (Mnemosyne), physical deletion creates unsolvable race conditions, stale read anomalies, and ambiguity regarding whether an absent cache record was deleted, evicted, or never created.
\end{enumerate}

In clinical and administrative domains, when an entity is retired, deactivated, cancelled, completed, entered in error, or otherwise logically removed, this action constitutes a domain-significant lifecycle transition rather than data destruction.

Furthermore, lower-level distributed cache capabilities (such as Infinispan's version-aware conditional removal primitive \texttt{removeWithVersion()} characterized in laboratory evaluations) must not be conflated with governed resource deletion. Similarly, long-term data archival and regulatory purge have sometimes been improperly assumed to be standard in-band application persistence features.

A definitive architectural decision is required to establish Harmonia's platform invariant regarding governed persisted information, the write semantics for authoritative Clinical operations (Task 08), the boundary between cache eviction and resource state, and the formal scope of archival and purge.

\subsubsection{Decision}

\begin{quote}
    \textbf{Harmonia does not provide physical DELETE semantics for governed persisted information. Logical deletion is a domain-appropriate lifecycle transition and therefore an authoritative UPDATE.}
\end{quote}

Harmonia establishes the following normative platform guardrail:

\begin{quote}
    \textbf{Harmonia SHALL NOT expose physical deletion as a normal operation for governed persisted information. Logical deletion SHALL be represented as a domain-appropriate lifecycle transition and processed as a concurrency-controlled authoritative update. Mneme cache eviction SHALL NOT be interpreted as deletion of authoritative state. Archival, retention-based purge and physical disposal are outside the current Harmonia framework scope.}
\end{quote}

\subsubsection{Four Core Distinctions}

To ensure conceptual clarity across architecture, design, and implementation, Harmonia strictly distinguishes four operational concepts:

\begin{enumerate}
    \item \textbf{Resource Lifecycle Transition}: An authoritative state change in Mnemosyne reflecting domain progression (illustrative examples include transitioning to \texttt{inactive}, \texttt{entered-in-error}, \texttt{cancelled}, \texttt{completed}, or \texttt{superseded}). A lifecycle transition is committed to Mnemosyne as a concurrency-controlled authoritative \texttt{UPDATE}, subject to version advancement, provenance capture, and Kleio audit logging.
    \item \textbf{Mneme Eviction}: The removal, invalidation, or expiration of a reconstructable active-state entry from the distributed cache (Mneme/Infinispan) according to cache capacity, TTL/idle policies, explicit invalidation, or cache-clearing operations. Cache eviction affects only the non-authoritative active cache working set; it does NOT alter, mutate, or delete authoritative durable state in Mnemosyne.
    \item \textbf{Archival}: The long-term migration, cold-storage management, or offline tiering of historical records. Archival is formally designated as outside the scope of the current Harmonia framework and must be managed by external operational and database lifecycle procedures.
    \item \textbf{Purge / Physical Disposal}: The physical destruction, cryptographic erasure, or permanent scrubbing of persisted records under statutory data-retention schedules or legal mandates. Retention-based purge and physical disposal are formally designated as outside the scope of the current Harmonia framework.
\end{enumerate}

\subsubsection{Write Semantics and Task 08 Consequence}

Harmonia categorizes authoritative write operations into distinct primitives:

\begin{itemize}
    \item \textbf{CREATE}: Establishes a new authoritative resource in Mnemosyne, establishes initial durable state and provenance, and participates in the version and coordination protocol defined by Task 08.
    \item \textbf{UPDATE}: Modifies existing authoritative resource attributes or lifecycle state in Mnemosyne, advances authoritative state under the version and concurrency control protocol defined by Task 08, and coordinates state with Mneme.
    \item \textbf{LIFECYCLE TRANSITION}: A semantic specialization of \textbf{UPDATE} that alters lifecycle attributes (illustrative examples include domain-appropriate fields such as FHIR \texttt{status}, \texttt{active}, or \texttt{verificationStatus}). It executes strictly through the authoritative \texttt{UPDATE} write pipeline.
    \item \textbf{DELETE (Physical)}: \textbf{Not supported.} Harmonia does not provide an authoritative physical delete path or database \texttt{DELETE} protocol for governed persisted information.
\end{itemize}

Implementation sequence Task 08 (``Establish authoritative Clinical writes'') is consequently scoped to design and implement authoritative \texttt{CREATE} and \texttt{UPDATE} write protocols between Mneme coordination and Mnemosyne persistence. Task 08 does not design, implement, or expose a physical authoritative \texttt{DELETE} path.

\subsubsection{Concurrency and Distributed Cache Semantics}

All lifecycle transitions execute through concurrency-controlled \texttt{UPDATE} protocols in Mnemosyne:

\begin{itemize}
    \item A lifecycle transition participates in the same concurrency-control and conflict-rejection mechanisms as any other authoritative update in Mnemosyne. Stale updates that conflict with committed durable state are rejected according to the concurrency protocol established in Task 08.
    \item Invalid lifecycle state transitions (such as attempting an illegal transition on an already terminal resource) are rejected by domain validation rules.
    \item In Mneme, Infinispan's \texttt{removeWithVersion()} is a version-aware conditional removal capability that Harmonia may use for reconstructable cache management and cluster coordination. It operates strictly within the non-authoritative active cache working set and does not define or execute governed resource deletion semantics.
\end{itemize}

\subsubsection{FHIR and Domain Lifecycle Semantics}

At the protocol and API surface:

\begin{itemize}
    \item External protocol interactions conceptually mapped to deletion (such as HTTP \texttt{DELETE} in FHIR REST APIs) do not execute physical database deletions in Mnemosyne.
    \item Instead, inbound protocol adapters and workflow pipelines translate such interactions into domain-appropriate lifecycle updates (illustrative examples include updating \texttt{active = false} or setting \texttt{status = 'entered-in-error'}), committing an authoritative update with associated provenance and audit evidence.
    \item Where domain or protocol specifications dictate specific responses (such as HTTP \texttt{410 Gone}) for retired or entered-in-error resources, the response is determined by evaluating governed lifecycle attributes in authoritative persistence, not by the physical absence of a record.
\end{itemize}

\subsubsection{Audit and Provenance}

In conformance with ADR-013 (Kleio Owns Audit Evidence and Provenance):

\begin{itemize}
    \item Accepted audit evidence in Kleio is append-only, permanent, and immutable.
    \item Every lifecycle transition produces immutable Kleio \texttt{AuditEvent} records capturing the actor, authorization context, transition timestamp, prior state, and resulting state.
    \item Because persisted entities remain permanently in Mnemosyne, audit trails and provenance graphs remain intact and verifiable throughout the system's life.
\end{itemize}

\subsubsection{Consequences}

\paragraph{Positive Consequences}

\begin{itemize}
    \item \textbf{Information Preservation}: Clinical and operational history is preserved without loss of longitudinal record integrity.
    \item \textbf{Referential Integrity}: Cross-resource references, patient links, and downstream integrations remain structurally sound without dangling foreign keys.
    \item \textbf{Audit Compliance}: Kleio audit evidence and provenance graphs remain anchored to permanent, verifiable historical records.
    \item \textbf{Concurrency Safety}: Lifecycle transitions are protected by concurrency control and version checks, preventing lost updates and race conditions.
    \item \textbf{Clear Architectural Boundaries}: Distinguishes non-authoritative Mneme cache eviction from authoritative Mnemosyne state transitions, and excludes complex archival/purge engines from platform scope.
\end{itemize}

\paragraph{Negative Consequences}

\begin{itemize}
    \item \textbf{Monotonic Storage Growth}: Persistent database tables grow monotonically, requiring operational database capacity monitoring and storage planning.
    \item \textbf{Query Filtering Discipline}: Search queries and read pipelines must consistently filter on domain lifecycle state (e.g. \texttt{active = true}, \texttt{status != 'entered-in-error'}) to prevent retired records from appearing in active clinical queries.
    \item \textbf{API Mapping Rigor}: Inbound protocol gateways must carefully map incoming deletion requests to domain-specific lifecycle state changes.
\end{itemize}

\subsubsection{Rejected Alternatives}

\begin{itemize}
    \item \textbf{Physical DELETE as Persistence Primitive}: Rejected because physical deletion destroys clinical history, invalidates foreign keys, breaks audit chains, and creates severe synchronization failures across distributed nodes.
    \item \textbf{Universal Generic Soft-Delete Flag (\texttt{is\_deleted = true})}: Rejected as the primary mechanism for domain lifecycle governance in favor of domain-native lifecycle attributes (e.g., FHIR \texttt{status}, \texttt{active}, \texttt{verificationStatus}, \texttt{period.end}). Generic boolean flags obscure domain semantics, fail to capture clinical context or reason for retirement, and complicate state validation. Any database-level columns (such as \texttt{is\_deleted}) are implementation artefacts and do not replace domain-level lifecycle representation.
    \item \textbf{In-Band Platform Purge and Archival Subsystem}: Rejected for the current Harmonia framework. Implementing compliant multi-jurisdictional data retention, cold-storage migration, and tombstone synchronization would add significant operational complexity not justified by the core integration mission.
\end{itemize}

\subsubsection{Relationship to Other Decisions}

This decision complements and reinforces:
\begin{itemize}
    \item \textbf{ADR-003 --- Mnemosyne Owns Durable Application State}
    \item \textbf{ADR-006 --- FHIR R5 Is the Clinical Interoperability Boundary}
    \item \textbf{ADR-010 --- Clinical Search Must Be Authoritative-Backed}
    \item \textbf{ADR-013 --- Kleio Owns Audit Evidence and Provenance}
    \item \textbf{ADR-018 --- Mnemosyne Defines the Authoritative Durable State Boundary}
    \item \textbf{ADR-019 --- Mneme Owns Distributed Resource Access and Coordination}
\end{itemize}

\subsubsection{Architectural Principle}

\begin{quote}
    \textbf{Governed information progresses through lifecycle states; it is never erased.}\\
    \textbf{State changes are authoritative updates; cache evictions are not deletions; archival and purge belong to external governance.}
\end{quote}

